% ============================================================================
%  3.7 VALIDACIONES Y MANEJO DE ERRORES
% ============================================================================

\section{Validaciones y manejo de errores}

\subsection{Validaciones implementadas}

El m\'odulo incluye las siguientes validaciones:

\begin{enumerate}
  \item \textbf{Autenticaci\'on:} Antes de inscribirse, se verifica
    que el usuario est\'a logueado. Si no lo est\'a, se muestra un
    alerta ``Iniciar sesi\'on''.

  \item \textbf{Duplicados:} La funci\'on \texttt{registerForEvent}
    verifica si ya existe una inscripci\'on activa antes de insertar.
    Si existe, lanza el error ``Ya est\'as inscrito en este evento''.

  \item \textbf{Restricci\'on SQL:} La tabla tiene una restricci\'on
    \texttt{UNIQUE (user\_id, event\_id)} que previene duplicados
    a nivel de base de datos.

  \item \textbf{Estado:} Solo se permiten los valores ``confirmed'',
    ``cancelled'' y ``waitlist'' mediante la restricci\'on CHECK.

  \item \textbf{Integridad referencial:} Las foreign keys
    \texttt{user\_id $\to$ profiles} y \texttt{event\_id $\to$ events}
    garantizan que las referencias sean v\'alidas.
\end{enumerate}

\subsection{Manejo de errores}

\subsubsection*{Errores de red / Supabase}

Todas las funciones de API capturan errores de Supabase y los
re-lanzan como \texttt{Error} con el mensaje descriptivo:

\begin{verbatim}
if (error) throw new Error(error.message);
\end{verbatim}

\subsubsection*{Errores en la interfaz}

Los handlers en React capturan los errores y muestran alertas:

\begin{verbatim}
const handleRegister = async (eventId) => {
  if (!user?.id) {
    Alert.alert('Iniciar sesión',
      'Debes iniciar sesión para inscribirte.');
    return;
  }
  try {
    await registerEvent({ userId: user.id, eventId });
  } catch (err) {
    const msg = err instanceof Error
      ? err.message
      : 'Error al inscribirte';
    Alert.alert('Error', msg);
  }
};
\end{verbatim}

\subsubsection*{Confirmaci\'on antes de cancelar}

Para la cancelaci\'on se muestra un di\'alogo de confirmaci\'on
antes de ejecutar la acci\'on, tanto en web como en m\'ovil:

\begin{itemize}
  \item \textbf{Web:} \texttt{window.confirm()}
  \item \textbf{M\'ovil:} \texttt{Alert.alert()} con opciones
    ``No'' y ``S\'i, cancelar''
\end{itemize}

\subsection{Tabla de validaciones}

\begin{table}[H]
\centering
\begin{tabular}{@{}lp{4cm}p{4cm}@{}}
\toprule
\textbf{Caso} & \textbf{Validaci\'on} & \textbf{Resultado} \\
\midrule
Usuario no logueado & \texttt{user?.id} null & Alerta de login \\
Ya inscrito & \texttt{checkRegistration} retorna datos & Error
  ``Ya est\'as inscrito'' \\
Datos inv\'alidos & UNIQUE constraint en DB & Error de Supabase \\
Red ca\'ida & Supabase lanza error & Alerta ``Error al
  inscribirte'' \\
Cancelar sin confirmar & Di\'alogo de confirmaci\'on & Acci\'on
  cancelada \\
\bottomrule
\end{tabular}
\caption{Tabla de validaciones del m\'odulo}
\end{table}
